\clearpage
\section{Design}

The design preserves the \textbf{domain structure} of assignment 3 and adds only
the distributed mechanisms that are necessary for clustering.

\subsection{Role-based node architecture}

Each JVM starts exactly one \textbf{cluster role}:
\begin{itemize}
    \item \texttt{control-unit}
    \item \texttt{keypad}
    \item \texttt{sensor}
\end{itemize}

This is more explicit than a monolithic startup. The node role determines
which local actor is started and how that node participates in the cluster.

\subsection{Single logical control unit}

The central design decision is that the \textbf{alarm state machine} still has
one logical owner: \texttt{ControlUnitActor}. In the clustered version it is not
spawned as an ordinary local actor, but exposed through a \textbf{cluster
singleton}. This guarantees that, even if multiple control-unit nodes are
started, only \textbf{one active control entity} manages the alarm state at a time.

\subsection{Discovery of distributed actors}

In assignment 3 the actors could receive direct references during local
construction. In assignment 4 this is no longer sufficient, because keypad and
sensor nodes may run on different JVMs.

The adopted solution is \textbf{receptionist-based discovery}:
\begin{itemize}
    \item the control unit registers itself as a discoverable service;
    \item keypads and sensors subscribe to that service and keep the discovered reference updated;
    \item sirens register themselves, and the control unit subscribes to the available siren instances.
\end{itemize}

\subsection{Recovery semantics}

The main behavioral extension with respect to assignment 3 is the state
\texttt{STARTUP\_RECOVERY}. After creation or recreation, the control unit:
\begin{itemize}
    \item does \textbf{not} assume the system is armed;
    \item does \textbf{not} assume the system is disarmed;
    \item ignores or logs sensor activations;
    \item accepts only a \textbf{correct PIN} to return to \texttt{DISARMED}.
\end{itemize}

This matches the specification exactly: state is intentionally \textbf{not
persisted}, so recovery must be \textbf{safe} rather than optimistic.

\subsection{Trade-offs}

This solution prefers \textbf{clarity} over maximum generality. It does not add
persistence, replicated state, or advanced failover policies, because those
features are outside the assignment goals. The chosen design is sufficient to
show:
\begin{itemize}
    \item distributed startup;
    \item remote discovery;
    \item single ownership of state;
    \item message-only interaction;
    \item safe recovery after recreation.
\end{itemize}
